\documentclass[10pt]{article}
\usepackage[T1]{fontenc}
\usepackage[utf8]{inputenc}
\usepackage[french]{babel}
\usepackage{amsmath,amssymb,mathrsfs,mathpazo}
\usepackage{geometry}
\geometry{a4paper,margin=2cm}
\usepackage{graphicx,array,tabularx,tabularray,colortbl,multirow,multicol,booktabs,enumitem,pifont,ulem,fancybox,diagbox,cancel}
\usepackage{pstricks,pstricks-add,pst-plot,pst-tree,pst-node,pst-text,pst-func,pst-math,pst-bspline,pst-3dplot,pst-eucl,pst-solides3d}
\let\pscalculate\undefined
\usepackage{pst-bezier}
\usepackage{tikz,pgfplots,tkz-tab}
\usetikzlibrary{arrows,arrows.meta,patterns,calc,positioning,trees,decorations.markings,shapes.geometric,shapes.misc,shapes.arrows,intersections,decorations.pathmorphing,decorations.pathreplacing,backgrounds,shadows,babel}
\pgfplotsset{compat=1.18}
\usepgfplotslibrary{fillbetween}
\usepackage[np]{numprint}
\usepackage{listings,makecell,pas-tableur,twoopt}
\let\barre\undefined
\usepackage[pstricks]{tablvar}
\definecolor{darkgreen}{RGB}{0,100,0}
\providecommand{\red}{\color{red}}
\providecommand{\blue}{\color{blue}}
\providecommand{\green}{\color{green}}
\providecommand{\black}{\color{black}}
\providecommand{\white}{\color{white}}
\providecommand{\gray}{\color{gray}}
\providecommand{\cyan}{\color{cyan}}
\providecommand{\magenta}{\color{magenta}}
\providecommand{\yellow}{\color{yellow}}
\providecommand{\violet}{\color{violet}}
\providecommand{\orange}{\color{orange}}
\providecommand{\euro}{\texteuro}
\providecommand{\EUR}{\texteuro}
\providecommand{\ieme}{\textsuperscript{e}}
\providecommand{\degres}{\ensuremath{^\circ}}
\providecommand{\degree}{\ensuremath{^\circ}}
\providecommand{\vv}[1]{\overrightarrow{#1}}
\providecommand{\Coord}[1]{#1}
\providecommand{\decosix}{}
\providecommand{\starredbullet}{\ensuremath{\bullet}}
\pagestyle{empty}
\setlength{\parindent}{0pt}
\usepackage{adjustbox}\DeclareUnicodeCharacter{2236}{:}\DeclareUnicodeCharacter{2190}{\ensuremath{\leftarrow}}
\providecommand{\R}{\mathbb{R}}
\providecommand{\N}{\mathbb{N}}
\providecommand{\Z}{\mathbb{Z}}
\providecommand{\Q}{\mathbb{Q}}
\providecommand{\D}{\mathbb{D}}
\providecommand{\C}{\mathbb{C}}
\providecommand{\e}{{\mathrm{e}}}
\providecommand{\vect}[1]{\overrightarrow{#1}}
\providecommand{\vectt}[1]{\overrightarrow{\mathrm{#1}}}
\providecommand{\vvt}[1]{\overrightarrow{\mathrm{#1}}}
\providecommand{\np}[2][]{{#2~#1}}
\providecommand{\vv}[1]{{\overrightarrow{#1}}}
\renewcommand{\d}{{\mathrm{d}}}
\providecommand{\degres}{{^\circ}}
\providecommand{\degre}{{^\circ}}
\providecommand{\blue}{\color{blue}}
\providecommand{\red}{\color{red}}
\providecommand{\enskip}{\quad}
\providecommand{\llbracket}{\lbrack\!\lbrack}
\providecommand{\rrbracket}{\rbrack\!\rbrack}
\providecommand{\euro}{{\text{€}}}
\providecommand{\no}{{\text{n°}}}
\providecommand{\interval}[2]{{[#1;#2]}}
\providecommand{\intervFO}[2]{{[#1;#2[}}
\providecommand{\dotfill}{\ldots}

\providecommand{\ds}{\displaystyle}
\providecommand{\textding}[1]{#1}
\providecommand{\Oij}{\(\mathrm{O};\vec{\imath},\vec{\jmath}\)}
\providecommand{\Oijk}{\(\mathrm{O};\vec{\imath},\vec{\jmath},\vec{k}\)}
\providecommand{\Ouv}{\(\mathrm{O};\vec{u},\vec{v}\)}
\providecommand{\intervFF}[2]{\ensuremath{[#1;#2]}}
\renewcommand{\euro}{\ifmmode\text{\texteuro}\else\texteuro\fi}
\DeclareUnicodeCharacter{00B0}{\ifmmode^{\circ}\else\textdegree\fi}
\DeclareUnicodeCharacter{20AC}{\ifmmode\text{\texteuro}\else\texteuro\fi}
% Énoncés : Éducation nationale. Transcription et corrigés : APMEP, Denis Vergès et contributeurs. Découpage et classement : CoopMaths, Sébastien Lozano, Éric Elter et contributeurs. Adaptation de la présentation : Maths Explorer.
% CC BY-SA : https://coopmaths.fr/www/about/
% https://coopmaths.fr/alea/static/e3c/2020/tex/e3c_2020_00_sujet15_3_cor.tex
\begin{document}
\noindent\textit{APMEP / CoopMaths. CC BY-SA. Présentation adaptée par Maths Explorer.}\par\medskip
\subsection*{1.}
	
	Avec \(\overrightarrow{AB} \left( \begin{array}{c} 12 \\ 6 \end{array} \right)\) et \(\overrightarrow{HC} \left( \begin{array}{c} 6 \\ -12 \end{array} \right)\),\\
	 on a \(\overrightarrow{AB} \cdot \overrightarrow{HC} = 12 \times 6 + 6 \times (-12) = 72 - 72 = 0\).
	
	De même avec \(\overrightarrow{AC} \left( \begin{array}{c} 12 \\ -12 \end{array} \right)\) et \(\overrightarrow{HB} \left( \begin{array}{c} 6 \\ 6 \end{array} \right)\), 
	\\on a \(\overrightarrow{AC} \cdot \overrightarrow{HB} = 12 \times 6 + 6 \times (-12) = 72 - 72 = 0\).
	
	\subsection*{2.}
	
	On a donc (CH) est perpendiculaire à (AB) et (BH) est perpendiculaire à (AC). Les droites (CH) et (BH) sont donc deux hauteurs du triangle ABC : elles sont donc sécantes en H orthocentre du triangle et la troisième hauteur est la droite (CH).
	
	\subsection*{3.}
	
	On a \(KA^2 = 9^2 + (-3)^2 = 81 + 9 = 90\),\\
	\(KB^2 = 3^2 + 9^2 = 9 + 81 = 90\),\\
	\(KC^2 = (-3)^2 + (-9)^2 = 9 + 81 = 90\).
	
	Or \(KA^2 = KB^2 = KC^2 = 90\) entraîne \(KA = KB = KC = R\). \\
	Le point K est équidistant de A, B et C : c’est donc le centre du cercle circonscrit au triangle ABC.
	
	\subsection*{4.}
	
	Avec \(M(8, 7)\) et avec \(G(g, g')\),
	\[
	\overrightarrow{AM} \left( \begin{array}{c} 12 \\ -3 \end{array} \right), \quad \overrightarrow{AM} \left( \begin{array}{c} 8 \\ -2 \end{array} \right).
	\]
	
	On a donc
	\[
	\overrightarrow{AG} \left( \begin{array}{c} 8 \\ -2 \end{array} \right) = \overrightarrow{AG} \left( \begin{array}{c} g + 4 \\ g' - 10 \end{array} \right).
	\]
	
	On en déduit que \(g = 8 - 4 = 4\) et \(g' = 10 - 2 = 8\). Donc \(G(4, 8)\).
	
	\subsection*{5.}
	
	On a
	\[
	\overrightarrow{GH} \left( \begin{array}{c} -2 \\ 2 \end{array} \right) \quad et \quad \overrightarrow{GK} \left( \begin{array}{c} 1 \\ 1 \end{array} \right).
	\]
	
	On a donc de façon évidente \(\overrightarrow{GH} = 2 \overrightarrow{GK}\) : les vecteurs sont colinéaires, les droites (GH) et (GK) sont parallèles mais ont le point G commun, donc les points G, H et K sont alignés.

\end{document}
